\section{Educational Implications}
\label{sec:education}

RQ4 asks what these results can teach; the answer is a designed resource, not a
measured outcome. No learners were observed or surveyed, so these statements describe what a lesson
makes visible, not what anyone learned.

The lesson has five observable objectives: explain why a model that scores well
historically may fail tomorrow; identify leakage in a code fragment; interpret a
skill score including a negative one; predict which families are scale-sensitive
and justify it mechanistically; and distinguish forecast error from profitability.
Each is anchored to a real artefact: the negative random-forest entry in
Table~\ref{tab:results} is the teaching moment for the third, and
Figure~\ref{fig:ablation} asks why the tree bars are exactly zero.

Its most valuable property is that the headline result is negative. A demonstration
where the sophisticated model wins teaches that complexity is good; this one
teaches the transferable question \emph{what is the baseline, and did we beat it?}
Five formative tasks with worked answers accompany it, and the notebook lets a
non-specialist re-run the comparison.
